\documentclass[fontset=none]{ctexart}

\setCJKmainfont{Noto Serif CJK SC}

\begin{document}

\section{绪论}
\subsection{研究背景}
操作系统是现代计算机体系结构的核心组成部分，承担着计算资源管理与上层软件运行支持的关键职能。随着软硬件技术的持续演进，传统操作系统的设计理念与实现机制正在面临新的需求与挑战。Jia等人（2025）在《Agent-Centric Operating System: A Comprehensive Review and Outlook for Operating System》中指出，随着 CPU、GPU 计算能力的提升，以及专为 AI 计算设计的 NPU 的出现与广泛应用，异构计算架构（Heterogeneous Computing Architecture）正逐渐成为主流。ARM big.LITTLE、Intel 混合核心以及 CPU–GPU 协同架构的提出，使系统在性能与能耗之间获得更灵活的平衡。然而，主流操作系统仍多基于同质多核（Symmetric Multi-Processing, SMP）假设：默认各核心性能相似、任务可在任意核心间自由迁移而不影响整体性能。这一假设在异构环境中显然失效。

在异构体系下，调度问题的数学结构已由传统的静态、线性、可分解形式，演化为动态、非线性、非对称、耦合且高维的复杂决策问题。传统算法依赖全局可知、确定性模型，而异构系统的运行状态不断变化、资源性能差异显著、目标函数多维且交叉耦合，使得启发式方法难以维持稳定最优解。如何在多类型处理器协同条件下实现自适应、可扩展的资源管理，成为操作系统设计的关键挑战，而线程调度正是其中最核心的问题之一。

现有操作系统的调度机制普遍建立在 SMP 模型之上，假设各计算核心在性能与能耗方面等价。这种假设在经典与现代系统中均有体现：教学操作系统 xv6 的调度循环在每个 CPU 核心上独立运行，所有核心采用相同策略；Zephyr 与 Linux 内核的 CFS（Completely Fair Scheduler）均基于 per-CPU 运行队列和负载均衡机制以实现“公平”分配，默认核心性能一致。FreeRTOS 在文档中虽提及可运行于多核场景（AMP 或 SMP），但其标准实现与示例主要针对单核或对称多核平台，未涉及异构核性能差异或异构感知调度策略。这些系统的调度算法——如先来先服务（FCFS）、最短作业优先（SJF）、时间片轮转（RR）——在设计时均假设任务可独立、核心可对称、环境可静态。

然而在异构体系下，性能差异、能耗约束、任务依赖与迁移成本交织，调度问题呈现出高度耦合的动态性。
传统调度机制因此暴露出三类局限：
1. 缺乏动态反馈：策略在运行期无法根据系统状态自我调整；
2. 局部最优陷阱：多核间负载均衡基于局部信息，难以形成全局最优分配；
3. 规则僵化：调度逻辑以固定启发式为核心，无法通过经验积累改进策略。

随着异构平台普及，这些假设的崩塌直接推动了从启发式调度向自学习调度的转变。
\subsection{国内外研究现状}
\subsection{研究内容与意义}
\end{document}
